\chapter{P\&L Analytics and the Strategy Lifecycle}\label{ch:hf:pnl-analytics-and-the-strategy-lifecycle}

A market-making strategy earned less this month than last. The spread was narrower, volume was higher, a competitor arrived on one venue, and
someone changed a parameter on the twelfth; the morning report has to say which of these did it. In this chapter's simulated year, month 9 made
\$45\,600 less than month 8: the narrower spread cost \$51\,500, higher volume gave back \$48\,000, the competitor took \$57\,400 and the parameter change
\$8\,100, with \$23\,400 of noise the other way. Estimated from the strategy's own data, without knowing what was planted, the competitor's share comes
out at \$53\,500 and the parameter's at \$8\,100.

\section{Daily attribution of a market-making book}

A market maker's daily P\&L (One Quant Book 8, chapter 1) splits into what it earned from the spread, what it lost to informed flow, what its
inventory made or lost, and its fees and rebates (One Quant Book 7, chapter 23, and chapter 1 here). For each venue and instrument, the expected part
is a product:
\[
\text{P\&L}=\underbrace{V}_{\text{market volume}}\times\underbrace{s}_{\text{our share}}\times\bigl(\underbrace{\kappa\,h}_{\text{capture}}
-\underbrace{a}_{\text{adverse selection}}\bigr),
\]
with $h$ the market half-spread, $\kappa$ the share of it the maker keeps (its capture ratio, net of fees and rebates) and $a$ the adverse selection
per share, measured by mark-outs. The factors are observable separately: $V$ and $h$ from the market, $s$ from the maker's fills, $\kappa$ and $a$ from
its fill prices and their mark-outs. A change in the product is a change in some factor, and the attribution's job is to say which.

The chapter's year (\cref{fig:hf:pnl-analytics-and-the-strategy-lifecycle:daily}): a maker on two venues, 10\% of 20 million shares a day on one and
12\% of 15 million on the other, keeping 80\% of a one-cent half-spread and losing 0.35 cent a share to adverse selection. Planted in it: the half-spread
narrows 10\% and market volume rises 15\% in month 9; a competitor arrives on venue B on day 175 and takes 30\% of the maker's share there and 15\%
of its capture ratio; a parameter change, from day 182, raises the maker's share by 10\% and its adverse selection by 0.10 cent a share on the days
it is switched on.

\begin{omfigure}
\begin{tikzpicture}
\begin{axis}[width=11cm,height=5.2cm,xlabel={\small trading day},ylabel={\small daily P\&L (\$ thousand)},xmin=0,xmax=252,ymin=0,ymax=30,
  tick label style={font=\footnotesize},grid=major,grid style={gray!25},table/col sep=comma,unbounded coords=jump,
  legend style={font=\scriptsize,at={(0.03,0.05)},anchor=south west},legend columns=2]
\fill[gray!12] (axis cs:168,0) rectangle (axis cs:189,30);
\addplot[gray,thin] table[x=day,y=pnl]{figdata/hft/28-pnl-analytics-and-the-strategy-lifecycle/daily.csv};
\addplot[omThm,thick] table[x=day,y=roll]{figdata/hft/28-pnl-analytics-and-the-strategy-lifecycle/daily.csv};
\draw[omDef,dashed] (axis cs:175,0) -- (axis cs:175,30);
\draw[omDef,dotted,thick] (axis cs:182,0) -- (axis cs:182,30);
\node[font=\scriptsize,anchor=north east,omDef] at (axis cs:174,29) {competitor};
\node[font=\scriptsize,anchor=north west,omDef] at (axis cs:183,29) {parameter};
\legend{daily,20-day mean}
\end{axis}
\end{tikzpicture}

\omcaption{A simulated year of a two-venue market maker's daily P\&L and its 20-day mean; shaded, month 9, when the half-spread narrows and market
volume rises; dashed, the competitor's arrival on venue B (day 175); dotted, the parameter change, switched on from day 182 on randomised days. Data:
\texttt{hf\_attrib.daily}.}\label{fig:hf:pnl-analytics-and-the-strategy-lifecycle:daily}
\end{omfigure}

The month-on-month change is attributed by changing the factors one at a time, in a fixed order, from month 8's values to month 9's
(\cref{lst:hf:pnl-analytics-and-the-strategy-lifecycle:change}): the half-spread, then market volume, then the competitor, then the parameter. The
spread and volume are measured from market data; the competitor's effect and the parameter's come from the next two sections. What is left is noise
and the approximation of using monthly averages.

\section{Parameter changes as experiments}

\begin{definition}[Parameter change control]\label{def:hf:pnl-analytics-and-the-strategy-lifecycle:control}
\emph{Parameter change control}\index{parameter change control} is the discipline of treating every change to a live strategy's parameters as a
recorded, reviewed and reversible experiment: who changed what, when and why, with the change applied to a randomised part of the strategy's activity
until its effect is measured.
\end{definition}

A parameter changed on the twelfth for everything from then on is confounded with everything else that happened after the twelfth. Applied on
randomised days instead (One Quant Book 7, chapter 21; Book 7's \texttt{firm.abtest} assigns them by a hash), its effect is estimated by comparing
treated and untreated days in the same weeks. Over the 70 days from day 182, 33 were treated: the parameter lowered daily P\&L by \$2\,670 (standard
error \$1\,070); adjusting for each day's market volume (CUPED) gives \$2\,550 (\$1\,070). The factors are sharper than the P\&L: the treated days' share
was exactly 10\% higher and their adverse selection 0.10 cent a share higher, because both are measured on every fill. The parameter buys volume at a
price in toxicity that exceeds its value; it should be rolled back.

\section{Capture decay and competition}

\begin{definition}[Capture decay]\label{def:hf:pnl-analytics-and-the-strategy-lifecycle:decay}
\emph{Capture decay}\index{capture decay} is the decline over time of the share of the spread a market-making strategy keeps, or of the share of the
flow it trades, as competitors copy it, venues change or the flow it served goes elsewhere.
\end{definition}

A competitor shows up first in the maker's share and capture ratio on the venue it joins. A one-sided CUSUM test (One Quant Book 7, chapter 13) on
venue B's daily share, against its target of 12\% with a slack of half a percentage point, crosses its threshold on day 176, one day after the arrival
(\cref{fig:hf:pnl-analytics-and-the-strategy-lifecycle:cusum}); the same test on the capture ratio crosses the same day. Comparing the days before
the change point with the untreated days after it estimates the competitor's effect: 29.9\% of the share and 14.9\% of the capture ratio, against the
planted 30\% and 15\%.

\begin{omfigure}
\begin{tikzpicture}
\begin{axis}[width=11cm,height=4.8cm,axis y line*=right,axis x line=none,xmin=150,xmax=199,ymin=0,ymax=1.2,
  ylabel={\small CUSUM statistic},tick label style={font=\footnotesize},table/col sep=comma]
\addplot[omDef,thick,dashed] table[x=day,y=cusum]{figdata/hft/28-pnl-analytics-and-the-strategy-lifecycle/cusum.csv};
\label{plot:hf:cusumstat}
\draw[gray,dotted] (axis cs:150,0.05) -- (axis cs:199,0.05);
\end{axis}
\begin{axis}[width=11cm,height=4.8cm,axis y line*=left,xmin=150,xmax=199,xlabel={\small trading day},ylabel={\small share on venue B (\%)},
  ymin=0,ymax=15,tick label style={font=\footnotesize},grid=major,grid style={gray!25},table/col sep=comma,
  legend style={font=\scriptsize,at={(0.03,0.05)},anchor=south west}]
\addplot[omThm,thick,const plot] table[x=day,y=share]{figdata/hft/28-pnl-analytics-and-the-strategy-lifecycle/cusum.csv};
\addlegendentry{share on venue B (left)}
\addlegendimage{/pgfplots/refstyle=plot:hf:cusumstat}\addlegendentry{CUSUM statistic (right)}
\end{axis}
\end{tikzpicture}

\omcaption{Detecting the competitor: the maker's daily share of venue B's volume and the one-sided CUSUM statistic of its shortfall from 12\% (slack
0.5 point, threshold 0.05, dotted), days 150 to 199; the share falls on day 175 and the statistic crosses on day 176; the treated days after day 182
show the parameter's 10\% more share. Data: \texttt{hf\_attrib.cusum\_path}.}\label{fig:hf:pnl-analytics-and-the-strategy-lifecycle:cusum}
\end{omfigure}

With both estimates, month 9 attributes as in \cref{fig:hf:pnl-analytics-and-the-strategy-lifecycle:month9}: of the \$45\,600 fall, the estimated
attribution gives the spread $-\$51\,500$, volume $+\$48\,000$, the competitor $-\$53\,500$ (planted: $-\$57\,400$) and the parameter $-\$8\,100$ (planted:
$-\$8\,100$); the residual, \$19\,600, is the month's noise in inventory, volume and spreads. The estimated competitor effect is smaller than the
planted one because the change point comes a day late.

\begin{omfigure}
\begin{tikzpicture}
\begin{axis}[width=11cm,height=5cm,ybar,bar width=10pt,xmin=-0.6,xmax=4.6,xtick={0,1,2,3,4},
  xticklabels={spread,volume,competitor,parameter,residual},ylabel={\small month 9 less month 8 (\$ thousand)},ymin=-70,ymax=60,
  tick label style={font=\footnotesize},grid=major,grid style={gray!25},table/col sep=comma,
  legend style={font=\scriptsize,at={(0.97,0.03)},anchor=south east}]
\addplot[fill=gray!40,draw=gray] table[x=i,y=truth]{figdata/hft/28-pnl-analytics-and-the-strategy-lifecycle/month9.csv};
\addplot[fill=omThm!60,draw=omThm] table[x=i,y=estimate]{figdata/hft/28-pnl-analytics-and-the-strategy-lifecycle/month9.csv};
\draw[black] (axis cs:-0.6,0) -- (axis cs:4.6,0);
\legend{planted truth,estimated}
\end{axis}
\end{tikzpicture}

\omcaption{The month-on-month change in P\&L attributed to the half-spread, market volume, the competitor on venue B and the parameter change, with
the planted values and with the values estimated from the strategy's data (the change point and the experiment); the residual is what neither
explains. Data: \texttt{hf\_attrib.month9}.}\label{fig:hf:pnl-analytics-and-the-strategy-lifecycle:month9}
\end{omfigure}

The order of attribution matters: the competitor's effect evaluated at month 9's narrower spread is smaller than at month 8's. A fixed, documented
order, or the average over orders (a Shapley attribution), is part of the report's definition; so is reporting the residual rather than hiding it.
Menkveld's decomposition of one high-frequency market maker (chapter 1: \euro1.55 a trade on the spread net of fees, \euro0.68 lost on positions) is
the same idea at the level of a firm's whole book.

\section{When to retire a strategy}

\begin{definition}[Strategy lifecycle]\label{def:hf:pnl-analytics-and-the-strategy-lifecycle:lifecycle}
The \emph{strategy lifecycle}\index{strategy lifecycle} is the sequence a trading strategy goes through, from research, review and a staged launch,
through monitored production and changes run as experiments, to reduction and retirement when its measured edge no longer covers its costs, each
stage with its criteria written in advance.
\end{definition}

A retirement rule is a kill criterion (One Quant Book 7, chapter 1) for a whole strategy: stop when the edge, measured over a window long enough to
be reliable, falls below what the strategy costs to run (people, technology, capital). In the simulated year the maker earned \$16\,600 a day before
month 8; after the competitor and with the parameter half on, \$12\,200. A rule that retires the strategy when its 60-day mean falls below a
cost of \$12\,000 a day fires on day 239; at \$13\,000 a day, on day 221; at \$10\,000, never. The rule is a decision to make before the decay, not
after it: the cost is known in advance, the window sets how fast the rule can act and how often it errs.

\section{Tutorial: which of these did it}\label{tut:hf:pnl-analytics-and-the-strategy-lifecycle:tut}

\begin{tutorial}
\textbf{Goal.} Attribute a month's change in a market maker's P\&L to its causes, with the causes estimated from the strategy's own data.
\textbf{End state:} the four figures and the numbers of the text.
\begin{tutsteps}
\item \textbf{The year} with planted events (\texttt{firm.mmattrib.Year}) and its daily attribution.
\item \textbf{The competitor}: CUSUM on venue B's share and capture ratio, then the effect from before and after the change point.
\item \textbf{The parameter}: the randomised experiment (\texttt{firm.mmattrib.experiment} on Book 7's \texttt{firm.abtest}).
\omcode{code/firm/mmattrib/firm_mmattrib.py}{83}{96}{Treated against untreated days: the effect on daily P\&L, with CUPED, and on share and adverse selection.}\label{lst:hf:pnl-analytics-and-the-strategy-lifecycle:exp}
\item \textbf{The attribution}, factor by factor.
\omcode{code/firm/mmattrib/firm_mmattrib.py}{119}{135}{Change the half-spread, then volume, then the competitor, then the parameter; what is left is the residual.}\label{lst:hf:pnl-analytics-and-the-strategy-lifecycle:change}
\end{tutsteps}
\textbf{What to change next.} Attribute with a Shapley average over orders; add a third venue where the competitor also arrives later; apply the
experiment to half the instruments instead of half the days.
\end{tutorial}

\section{Build: the attribution engine}\label{bld:hf:pnl-analytics-and-the-strategy-lifecycle:mmattrib}

\begin{build}
\bfield{Purpose} Attribute a market maker's P\&L to its factors daily and month on month, detect changes in capture, run parameter changes as
experiments, and apply a retirement rule.
\bfield{Interface} \texttt{Year(seed)} and its planted-event parameters, \texttt{attribution(year)}, \texttt{cusum\_down(x, target, k, h)}, \texttt{experiment(year, first,
last)}, \texttt{month\_change(year, m0, m1, est)}, \texttt{retire\_day(pnl, window, floor)}. Built on Book 7's \texttt{firm.abtest}.
\bfield{Rules} Dollars and cents per share; a fixed order of attribution; the residual reported.
\bfield{Acceptance tests} \texttt{code/firm/mmattrib/tests/}: the parts add to the total and the observed capture ratio equals the planted one; the
competitor lowers venue B's share and the experiment starts on its day; CUSUM fires the day after a step and never without one; the experiment
recovers the planted share and adverse selection; without noise the attribution leaves under 5\% unexplained; the retirement rule fires after a
drop and not before.
\bfield{Stretch} Shapley attribution; attribution by instrument and venue at once; sequential tests for the experiment.
\end{build}

\begin{omsources}
\item A.~J.~Menkveld, High frequency trading and the new market makers, \emph{Journal of Financial Markets} 16(4), 2013, 712--740.
\item One Quant Book 7, chapters 1, 13 and 21 (kill criteria, CUSUM tests, experiments).
\end{omsources}

\section{Exercises}

\begin{exercise}[$\star$]\label{exo:hf:pnl-analytics-and-the-strategy-lifecycle:1}
A venue trades 15 million shares a day; the maker's share is 12\%, it keeps 80\% of a one-cent half-spread and loses 0.35 cent a share. What does it
make a day there?
\end{exercise}

\begin{exercise}[$\star$]\label{exo:hf:pnl-analytics-and-the-strategy-lifecycle:2}
The competitor cuts the share to 8.4\% and the capture ratio to 68\%. What does the venue make now, at the same spread?
\end{exercise}

\begin{exercise}[$\star$]\label{exo:hf:pnl-analytics-and-the-strategy-lifecycle:3}
The parameter raises the share by 10\% and adverse selection by 0.10 cent. Does it pay on the example venue before the competitor?
\end{exercise}

\begin{exercise}[$\star\star$]\label{exo:hf:pnl-analytics-and-the-strategy-lifecycle:4}
Why is a parameter changed for all days from the twelfth hard to evaluate?
\end{exercise}

\begin{exercise}[$\star\star$]\label{exo:hf:pnl-analytics-and-the-strategy-lifecycle:5}
Why does the order of attribution change the competitor's share of the fall?
\end{exercise}

\begin{exercise}[$\star\star$]\label{exo:hf:pnl-analytics-and-the-strategy-lifecycle:6}
Why are the share and adverse selection better measures of the parameter's effect than daily P\&L?
\end{exercise}

\begin{exercise}[$\star\star\star$]\label{exo:hf:pnl-analytics-and-the-strategy-lifecycle:7}
\emph{Coding.} Run the attribution on seeds 2 and 3. How large is the residual, and how close are the estimated competitor and parameter effects to
the planted ones?
\end{exercise}

\begin{exercise}[$\star\star\star$]\label{exo:hf:pnl-analytics-and-the-strategy-lifecycle:8}
\emph{Find the flaw.} ``P\&L fell after we changed the parameter on the twelfth, so the parameter change caused it; revert it.''
\end{exercise}

\section{Problem: Which of These Did It}

\begin{problem}[{Weekend problem --- which of these did it}]\label{pb:hf:pnl-analytics-and-the-strategy-lifecycle:1}
A market maker's morning report must explain why month 9 made less than month 8.

\medskip\noindent\textbf{Part I --- Attribution.}
\begin{enumerate}
\item Write the expected daily P\&L as a product of factors and say where each is measured.
\item Describe the simulated year and its planted events.
\item Describe the sequential attribution.
\item What does Menkveld's decomposition show at the level of a firm?
\end{enumerate}

\medskip\noindent\textbf{Part II --- Estimating the causes.}
\begin{enumerate}[resume]
\item Define \omterm{def:hf:pnl-analytics-and-the-strategy-lifecycle:control}{parameter change control}.
\item Give the experiment's estimates, with and without CUPED.
\item Define \omterm{def:hf:pnl-analytics-and-the-strategy-lifecycle:decay}{capture decay} and describe the CUSUM detection.
\item Give the estimated competitor effect against the planted one.
\end{enumerate}

\medskip\noindent\textbf{Part III --- The attribution.}
\begin{enumerate}[resume]
\item Give the planted and estimated attribution of month 9.
\item Why is the estimated competitor effect smaller?
\item What is the residual?
\item Why must the order of attribution be fixed and documented?
\end{enumerate}

\medskip\noindent\textbf{Part IV --- The verdict.}
\begin{enumerate}[resume]
\item State the \emph{named result}: the attribution of the month-on-month P\&L change to spread, volume, competition and the parameter change,
against the planted truth.
\item Should the parameter be kept?
\item Define the \omterm{def:hf:pnl-analytics-and-the-strategy-lifecycle:lifecycle}{strategy lifecycle}.
\item When does the retirement rule fire, and on what does it depend?
\item What should the report show every morning?
\item How would a Shapley attribution differ?
\item What would the competitor do to a strategy that cannot randomise its changes?
\item In one sentence: what does attribution make possible?
\end{enumerate}
\end{problem}

\section{Interview questions}

\begin{interviewq}[$\star$ \iqroles{trader}]\label{iq:hf:pnl-analytics-and-the-strategy-lifecycle:1}
Your strategy's P\&L halved this week. What do you look at first?
\end{interviewq}

\begin{interviewq}[$\star\star$ \iqroles{researcher}]\label{iq:hf:pnl-analytics-and-the-strategy-lifecycle:2}
How would you detect that a competitor has joined your venue from your own data?
\end{interviewq}

\begin{interviewq}[$\star\star$ \iqroles{developer}]\label{iq:hf:pnl-analytics-and-the-strategy-lifecycle:3}
Design the system that records every parameter change and its experiment assignment, and produces the attribution report by 7 a.m.
\end{interviewq}

\begin{interviewq}[$\star\star$ \iqroles{risk}]\label{iq:hf:pnl-analytics-and-the-strategy-lifecycle:4}
A trader wants to change a live parameter now, not as an experiment. What do you require?
\end{interviewq}

\begin{interviewq}[$\star\star$ \iqroles{researcher}]\label{iq:hf:pnl-analytics-and-the-strategy-lifecycle:5}
How many days does an experiment need to detect a \$2\,000-a-day effect with a daily standard deviation of \$4\,500?
\end{interviewq}

\begin{interviewq}[$\star\star\star$ \iqroles{researcher}]\label{iq:hf:pnl-analytics-and-the-strategy-lifecycle:6}
Show that a sequential attribution of a product of two factors assigns the interaction term to whichever factor is changed second.
\end{interviewq}
